\subsection*{\theappendixitem.\ Inductive limit of local rings}
\phantomsection
\addcontentsline{toc}{subsection}{\theappendixitem.\ Inductive limit of local rings}

Let I be a nonempty right-filtering preordered set, and let \((\mathbf{A}_{\alpha}, \varphi_{\beta\alpha})\) be an inductive system of rings relative to I. Suppose that, for every \(\alpha \in I\) , the ring \(A_{\alpha}\) is local, with maximal ideal \(m_{\alpha}\) , that the homomorphisms \(\varphi_{\beta\alpha}\) are local and flat, and that one has \(\varphi_{\beta\alpha}(m_{\alpha}) A_{\beta} = m_{\beta}\) for \(\beta \geqslant \alpha\) . Denote by A the inductive limit of the \(A_{\alpha}\) , and for every \(\alpha \in I\) , let \(\varphi_{\alpha}: A_{\alpha} \to A\) be the canonical homomorphism.

PROPOSITION 1. --- a) The ring A is local, with maximal ideal \(\mathfrak{m} = \varinjlim \mathfrak{m}_{\alpha}\) . For every \(\alpha \in \mathbf{I}\) , the homomorphism \(\varphi_{\alpha}\) is local and flat, and one has \(\varphi_{\alpha}(\mathfrak{m}_{\alpha})\) A = \(\mathfrak{m}\) .

\begin{enumerate}
\def\labelenumi{\alph{enumi})}
\setcounter{enumi}{1}
\item
  If \(\mathbf{A}_{\alpha}\) is Noetherian for every \(\alpha \in \mathbf{I}\) , then \(\mathbf{A}\) is Noetherian.
\item
  Set \(m = \varinjlim m_{\alpha}\) ; it is an ideal of A. The quotient ring A/m is the inductive limit of the fields \(A_{\alpha}/m_{\alpha}\) , it is therefore a field (A, I, p.~116, prop. 3). Moreover, every element of A - m is invertible in A: indeed, let \(x \in A - m\) ; there exists \(\alpha \in I\) and \(\xi \in A_{\alpha}\) such that \(x = \varphi_{\alpha}(\xi)\) ; one has \(\xi \notin m_{\alpha}\) , hence \(\xi\) is invertible in \(A_{\alpha}\) and x is invertible in A. Consequently, A is a local ring with maximal ideal m. Let \(\alpha \in I\) . From the relations \(\varphi_{\beta\alpha}(m_{\alpha}) A_{\beta} = m_{\beta}\) for \(\beta \geqslant \alpha\) , we deduce, by passing to the inductive limit, \(\varphi_{\alpha}(m_{\alpha}) A = m\) ; finally, the homomorphism \(\varphi_{\alpha}\) is flat by I, § 2, n° 7, prop. 9.
\item
  Let \(\hat{\mathbf{A}}\) be the separated completion of \(\mathbf{A}\) for the m-adic topology and \(\pi\) the canonical map from \(\mathbf{A}\) to \(\hat{\mathbf{A}}\) . Suppose the rings \(\mathbf{A}_{\alpha}\) are Noetherian. Fix \(\alpha \in \mathbf{I}\) and prove that the ring \(\hat{\mathbf{A}}\) is Noetherian and flat over \(\mathbf{A}_{\alpha}\) . By hypothesis, \(m_{\alpha}\) is a finitely generated ideal of \(\mathbf{A}_{\alpha}\) , hence \(m = \varphi_{\alpha}(m_{\alpha})\) . \(m\) is a finitely generated ideal of \(\mathbf{A}\) . It follows that the maximal ideal \(\hat{m}\) of \(\hat{\mathbf{A}}\) is equal to \(m\hat{\mathbf{A}}\) (III, § 2, n° 12, cor. 2 to prop. 16 and n° 13, prop. 19), hence is finitely generated. Consequently, the ring \(\hat{\mathbf{A}}\) is Noetherian (loc. cit., n° 10, cor. 5 to th. 2). On the other hand, for every \(n \in \mathbf{N}\) , the quotient \(\hat{\mathbf{A}}/\hat{m}^{n}\) is isomorphic to \(\mathbf{A}/m^{n}\) (loc. cit., n° 12, cor. 2 to prop. 16 and formula (21)), which means that \(\hat{\mathbf{A}}/\pi \circ \varphi_{\alpha}(m_{\alpha}^{n})\) \(\hat{\mathbf{A}}\) is isomorphic to \(\mathbf{A} \otimes_{\mathbf{A}_{\alpha}} (\mathbf{A}_{\alpha}/m_{\alpha}^{n})\) ; since \(\mathbf{A}\) is a flat \(\mathbf{A}_{\alpha}\)-module, the \((\mathbf{A}_{\alpha}/m_{\alpha}^{n})\) -module \(\hat{\mathbf{A}}/\pi \circ \varphi_{\alpha}(m_{\alpha}^{n})\) \(\hat{\mathbf{A}}\) is flat for every \(n \in \mathbf{N}\) . By III, § 5, n° 4, prop. 2, the \(\mathbf{A}_{\alpha}\) -module \(\hat{\mathbf{A}}\) is ideally separated for \(m_{\alpha}\) ; by loc. cit., n° 2, th. 1, the \(\mathbf{A}_{\alpha}\) -module \(\hat{\mathbf{A}}\) is therefore flat. It follows by passing to the inductive limit that \(\hat{\mathbf{A}}\) is (faithfully) flat over \(\mathbf{A}\) (I, § 2, n° 7, prop. 9), hence that \(\mathbf{A}\) is Noetherian (I, § 3, n° 5, corollary to prop. 8).
\end{enumerate}

\refstepcounter{appendixitem}
